\section{Выводы}

В лабораторной работе реализован поиск одного образца при помощи Z-функции над массивом слов.
Решение не сводит задачу к посимвольному поиску: один элемент сравнения -- это целое слово входного
алфавита. Регистронезависимость обеспечивается нормализацией всех слов перед построением объединённой
последовательности.

Ключевая идея алгоритма состоит в том, что образец и текст объединяются через разделитель, после чего
достаточно одного построения Z-функции. Значение Z-функции в позиции текстовой части сразу показывает,
сколько слов образца совпало с текстом с этой позиции. Если совпала вся длина образца, программа берёт
сохранённую позицию соответствующего слова и выводит ответ.

Бенчмарк подтвердил ожидаемое поведение. На случайных коротких сравнениях наивный поиск может быть
быстрее из-за малых констант и ранних несовпадений, однако в худшем случае его время быстро растёт:
для образца длины 50 и текста из \(50000\) одинаковых слов Z-поиск оказался примерно в
\(27.58\times\) быстрее. При масштабировании текста с \(100000\) до \(1000000\) слов время Z-поиска
выросло примерно в \(10.5\) раза, что согласуется с линейной оценкой.
\pagebreak
